\chapter{Dyson Rings and Solar Slingshots}
\label{ch:25}

The most extravagant of the concepts in the author's notes is the collection of stellar energy on a scale that dwarfs the whole of human use, and it is the concept with the longest history of misunderstanding, which begins with its name. Dyson's 1960 paper~\cite{dyson1960} proposed that an advanced civilization might be detected by the infrared radiation of structures that intercept the light of its star, and he proposed these structures as a swarm of independent bodies in orbit, a loose shell of many objects, and not a solid sphere. The solid sphere belongs to the literature of fiction. The notes use the terms ring and swarm, together with a gravitational slingshot around the Sun, and this chapter begins by distinguishing the architectures that the terms can denote, because the argument of the earlier chapters, that a proposal is incomplete unless its mechanism is stated, applies with special force to a concept of which a single name stands for physically very different objects. The closest experimentally grounded precedent for propulsion by light on a large sail is the analysis of laser-pushed lightsails~\cite{forward1984}, which concerns a different problem from the one considered here.

Consider first the rigid shell, a solid surface enclosing the star. The gravitational field inside a uniform spherical shell vanishes, a result of Newton, so that the shell exerts no force on the star and the star none on the shell, and the position of the shell relative to the star is in neutral equilibrium: any displacement leaves it displaced, and the slightest drift brings one side toward the star and the other away, with no restoring force. The compressive stresses in a shell of radius one astronomical unit under the star's gravity exceed the strength of any known material by many orders of magnitude, and the mass of the shell is enormous. Even a shell with an areal density of one kilogram per square meter, a film a millimeter thick of material of the density of water, has a mass of $2.8\times10^{23}$ kg, which is eighty-five per cent of the mass of Mercury. The rigid shell is excluded for reasons of stability, of strength, and of mass.

The rigid ring is a smaller structure with a similar defect. A solid ring in orbit around a central mass is in unstable equilibrium, as Maxwell showed in 1859~\cite{maxwell1859} in his essay on the stability of the motion of Saturn's rings: a displacement of the ring toward one side increases the attraction on that side, and the displacement grows. Maxwell concluded that the rings of Saturn must be composed of many small separate particles, and the conclusion has been confirmed by observation. A ring around a star can be held in place only by active control, with thrusters or by a continuing exchange of momentum, and a rigid ring of astronomical size is a structure that must be controlled at every moment, which compounds the difficulty. A third architecture is the statite shell, a surface of reflective sails held at rest by radiation pressure against gravity, and it has a precise criterion. A perfectly reflecting sail at any distance from the star experiences radiation force and gravitational force that both vary as the inverse square of distance, so the balance depends only on the areal density, and the critical value~\cite{mcinnes1999} is $\sigma_c=L/(2\pi cGM)=1.53\ \mathrm{g\,m^{-2}}$ for the Sun. A sail with a lower areal density is pushed outward, a sail with a higher one is pulled in, and a shell of such sails at one astronomical unit has a mass of $4.3\times10^{20}$ kg or more, about $0.13$ per cent of the mass of Mercury, if it is a complete shell at the critical density. The film must be thinner than a few micrometers of a material of ordinary density, and it must be reflective, thermally stable, and kept in position against the perturbations of the other bodies of the solar system.

The fourth architecture is the swarm, in which independent collectors in orbit around the star, each a sail, a mirror, or a platform with a radiator, are numerous enough to intercept a given fraction of the star's output without being rigidly connected. This is Dyson's own picture. It requires no structural strength beyond that of the individual collectors, since each is in free fall, and it requires control only for the avoidance of collision and the maintenance of orbit, tasks that scale with the number of collectors. A ring in the notes is best understood as a particular arrangement of the swarm, a band of collectors distributed around a circular orbit, whose width determines the fraction of the stellar output that it intercepts. For a band of width $W$ at orbital radius $r$ with its collectors oriented to face the star, the intercepted fraction is $W/(2r)$. A band of width one million kilometers at one astronomical unit intercepts $0.33$ per cent of the output, $1.3\times10^{24}$ W, and has an area of $9.4\times10^{20}\ \mathrm{m^2}$ and, at a mass per unit area of one gram per square meter, a mass of $9.4\times10^{17}$ kg. The numbers show two things. The mass needed to intercept a given fraction is small, a small fraction of the mass of the asteroid belt, because collectors can be thin. And the power intercepted is vastly greater than anything that could be used: it exceeds present human use by nearly eleven orders of magnitude. The sizes at which the concept begins to be relevant to human purposes are accordingly many orders below the ring of a million kilometers. A swarm intercepting a petawatt, about fifty-five times present use, requires a collecting area of $2.4\times10^{12}\ \mathrm{m^2}$ at an efficiency of thirty per cent and, at one gram per square meter, $2.4\times10^6$ tonnes, a modest mass by the standards of terrestrial industry, but a mass to be placed in space, which is the whole difficulty.

The difficulty determines the sequence, and the quantity that determines it is the fraction of the energy that is useful in space. The energy collected cannot all be sent to the surface, because, as Chapter 22 showed, the surface of the Earth can shed only some tens of terawatts of additional heat without a measurable change of climate, so the energy would have to be used at or near the collectors for purposes such as manufacture, propulsion, and computation in space. This sets the purpose of the swarm: it is not an electric utility for the Earth, and a proposal that presented it as such would founder on the heat budget. It is infrastructure for an economy in space, and its justification depends on the existence of that economy.

The material for the swarm would have to come from somewhere other than the Earth, and the notes propose the dismantling of Mercury, a planet of $3.3\times10^{23}$ kg whose gravitational binding energy is $1.8\times10^{30}$ J. This is the energy equivalent of $4{,}700$ seconds of the Sun's output, and to disassemble the planet within a century at perfect efficiency would require $5.7\times10^{20}$ W, which is $1.5\times10^{-6}$ of the solar luminosity. The figure is of the order given in the literature on the subject. It shows that the process is limited by the engineering of mining, transport, and processing and not by the energy, and it shows also the scale of the intervention, the removal of a planet, about which Chapter 13 and Chapter 16 require that authority and consent be addressed before the engineering is.

The figure also invites a correction of scale. Whole-planet disassembly is the extreme case, and the uses that the book has reason to consider need a vanishing fraction of Mercury. A collecting swarm of one petawatt needs about $2.4\times10^{9}$ kg, which is $7\times10^{-15}$ of the planet, and a heat shield or a fleet of battery factories needs less. The critics of the concept have been right about the planet and have not been addressing the uses. Knapp's objection of 2012~\cite{knapp2012} is that the energy and effort of dismantling Mercury would repay themselves only on a timescale far beyond any benefit, and that the same resources would do more good elsewhere. The objection holds for a scheme that spends energy at the outset and collects the return later, since a static ledger of this kind is dominated by the binding energy. It does not hold for a scheme that bootstraps, in which the collectors that the first mining produces power the next round of mining. A film of aluminium of areal density one gram per square meter has an embodied energy of about $0.2\ \mathrm{MJ\,m^{-2}}$ at $200\ \mathrm{MJ\,kg^{-1}}$. At the Earth's distance a collector at thirty per cent efficiency returns $408\ \mathrm{W\,m^{-2}}$ and repays the energy of its own film in about $490$ s, and at Mercury's orbit, where the solar flux is $9{,}083\ \mathrm{W\,m^{-2}}$, $6.7$ times the flux at the Earth, it returns $2.7\ \mathrm{kW\,m^{-2}}$ and repays it in about $73$ s. The film is not the whole cost, since each square meter of collector requires machinery to mine, refine, and deploy it. If the machinery is a hundred kilograms per square meter of collector at the same specific embodied energy, the payback is $1.55$ years at the Earth's distance and $0.23$ years at Mercury's, and the doubling times of an installation that reinvests all its output are $1.08$ and $0.16$ years respectively. Growth from one square kilometer to the $2.4\times10^{12}\ \mathrm{m^2}$ of a petawatt swarm requires about twenty-one doublings, which at Mercury is about three and a half years under these stipulated inputs. The bootstrapping answer is therefore valid, and it is a conditional one: the doubling time is set by the machinery's embodied energy and not by the Sun, the machinery must be able to reproduce itself from regolith with a closure that is not presently demonstrated, and a growth law of this kind is exactly the kind of process for which Chapter 16 requires a stop rule, because a process that doubles in a year is also a process that cannot be recalled after twenty doublings.

The gravitational slingshot around the Sun is a different kind of proposal: a means of sending a payload out of the solar system at a high speed. The principle is the Oberth effect~\cite{oberth1929}: a rocket burn made at high speed deep in a gravitational well changes the payload's kinetic energy by an amount proportional to its speed at the time of the burn, so that an increment of velocity delivered at perihelion is worth more than the same increment delivered far away. If the payload passes the Sun at a perihelion speed $v_p$ and burns to add $\Delta v$, the energy gained per unit mass is $v_p\Delta v+\tfrac12\Delta v^2$. The cost of the strategy is the dive: to bring a payload from the Earth's orbit to a perihelion of ten solar radii requires cancelling most of the Earth's orbital velocity, a velocity change of $20.9$ km/s, and the payload must survive the heat. At ten solar radii the solar flux is $6.3\times10^5\ \mathrm{W\,m^{-2}}$, which is $462$ times that at the Earth, and a blackbody surface in equilibrium with this flux, radiating from one side, would reach $1{,}825$ K. The Parker Solar Probe has come to within about ten solar radii, at a speed of about $190$ km/s, protected by a carbon-composite shield, and this is the heat load that a larger payload would face. The comparison with a direct departure is informative. A payload that is to leave with a hyperbolic excess speed of $47$ km/s needs $33.4$ km/s of velocity increment from Earth orbit if it goes directly, and $30.9$ km/s, composed of $20.9$ km/s for the dive and $10$ km/s at perihelion, if it uses the slingshot. For a smaller excess speed the direct departure is cheaper; at $15$ km/s the dive costs $25.9$ km/s and the direct departure $15.1$ km/s. The slingshot thus breaks even near an excess speed of forty kilometers per second and is advantageous only above it. It is a technique for very high speeds, and in the context of the book it would matter for missions that no program of present obligation requires. It is classified accordingly in the third readiness class, and its relation to the other orbital programs is that of a possible later extension.

Two uses of Mercury material bear on the slingshot and on the economy of the swarm, and both are proposals of the author's notes. The first is the heat shield. A payload that dives to ten solar radii must carry a shield, and the shield is a considerable fraction of its mass, as in the Parker Solar Probe, where the shield is of the order of a tenth of the spacecraft. If the shield is made from material mined and fabricated at Mercury, its mass need not be lifted from the Earth, and the dive can start from Mercury's orbit, where the orbital speed is $47.9$ km/s and the dive to ten solar radii costs $25.7$ km/s, against $20.9$ km/s from the Earth's orbit. The perihelion speed on the dive ellipse is $184.6$ km/s, and a burn of $15$, $20$, or $25$ km/s there gives hyperbolic excess speeds of $40.9$, $60.8$, and $75.9$ km/s at total increments of $40.7$, $45.7$, and $50.7$ km/s. A direct departure from Mercury's orbit to the same speeds costs $31.2$, $43.1$, and $53.8$ km/s, so that from Mercury the slingshot breaks even near $60$ km/s, which is a higher speed than the forty kilometers per second from the Earth. The result is not what the proposal expects. Mercury staging helps by supplying shield material and solar power on site, and the dive is no cheaper there, because the orbital speed to be cancelled is larger. The shield is thus an argument for building at Mercury and not an argument that the slingshot has become attractive, and the classification of the slingshot in the third readiness class is unchanged.

The second use is the remote production and charging of batteries in orbit about Mercury. The solar flux there is $6.7$ times that at the Earth, and the planet supplies the metals, silicon, and carbon of cells, so a battery factory powered by collectors in Mercury orbit is a plausible first customer of the swarm. The destination of the batteries sets the limit on the proposal. To move a battery from Mercury to the Earth by a Hohmann transfer costs $9.61$ km/s at departure and $7.53$ km/s at arrival, $17.1$ km/s in all, in a transit of $105$ days, and the kinetic energy that this increment represents, $\tfrac12\Delta v^2$, is $147\ \mathrm{MJ\,kg^{-1}}$ of payload. A battery of $250\ \mathrm{Wh\,kg^{-1}}$ holds $0.9\ \mathrm{MJ\,kg^{-1}}$, so that moving the battery by rocket costs about a hundred and sixty times the energy it stores. Even escape from Mercury alone, $4.25$ km/s, costs $9.0\ \mathrm{MJ\,kg^{-1}}$ for a mass driver, ten times the stored energy. The stored energy of a battery cannot, therefore, be exported by propulsion in any economic sense. The battery factories at Mercury have a use only if the energy is consumed where the batteries are: in tugs and miners that work in Mercury orbit, in the charging of vehicles that return only a short distance, or in swap stations in which a battery is exchanged between vehicles that already travel the route. Export of the energy itself as a beam is a different proposal with its own losses and hazards, and it faces the heat budget of Chapter 22. The two uses together give the result that the Mercury program is justified, if at all, as an industrial base for the orbital economy, which is the position the chapter has reached by the argument from scale.

\section*{Formalization}

Four results are stated: the statite criterion, the swarm geometry, the Oberth gain, and the cost of the dive.

For a flat perfectly reflecting sail of areal density $\sigma$ at distance $r$ from a star of luminosity $L$ and mass $M$, facing the star, the radiation force per unit area is $2L/(4\pi r^2c)$ and the gravitational force per unit area is $\sigma GM/r^2$.

\begin{proposition}
The sail is in force balance at every distance if and only if $\sigma=\sigma_c=L/(2\pi cGM)$. For $\sigma>\sigma_c$ the net force is toward the star and for $\sigma<\sigma_c$ away from it, at every $r$. For the Sun, $\sigma_c=1.53\ \mathrm{g\,m^{-2}}$ (or $0.77\ \mathrm{g\,m^{-2}}$ for a perfect absorber). A complete shell at radius $r$ at the critical density has mass $4\pi r^2\sigma_c$, which for $r=1$ AU is $4.3\times10^{20}$ kg.
\end{proposition}

\begin{proof}
Equating $L/(2\pi r^2c)$ to $\sigma GM/r^2$ gives $\sigma_c$, independent of $r$. The sign of the net force follows from the sign of $\sigma_c-\sigma$. The shell mass is area times density, with the area $4\pi(1.496\times10^{11})^2=2.81\times10^{23}\ \mathrm{m^2}$.
\end{proof}

For a rigid shell of areal density $1\ \mathrm{kg\,m^{-2}}$ the mass is $2.81\times10^{23}$ kg, which is $0.85$ of the mass of Mercury and exceeds the critical value for a statite by a factor of $650$, so such a shell cannot be supported by radiation.

For the swarm, let a band of collectors at orbital radius $r$ and width $W$ face the star, each with an effective collecting area per unit of band area $\alpha$, and let the optical-to-useful conversion efficiency be $\eta$.

\begin{proposition}
For $W\ll r$ the fraction of stellar output intercepted by the band is $f=W/(2r)$, the useful power is $\eta fL$, and the band area is $2\pi rW$, whose mass at areal density $\sigma_a$ is $2\pi rW\sigma_a=4\pi r^2f\sigma_a$. Hence the mass per unit intercepted power is $4\pi r^2\sigma_a/L$ per unit of fraction, independent of the geometry of the band, and equal to $\sigma_a/(\Phi\eta)$ per unit of useful power, where $\Phi=L/(4\pi r^2)$ is the flux at the band.
\end{proposition}

\begin{proof}
A band of width $W$ at radius $r$ presents a projected area $2\pi rW$ to the star, and the star's output passing through the sphere of radius $r$ is distributed over area $4\pi r^2$, giving the fraction $2\pi rW/(4\pi r^2)=W/(2r)$. The remaining statements follow from the definitions.
\end{proof}

With $L=3.828\times10^{26}$ W and $r=1.496\times10^{11}$ m, a band of $W=10^9$ m intercepts $f=3.34\times10^{-3}$, i.e. $1.28\times10^{24}$ W, has area $9.4\times10^{20}\ \mathrm{m^2}$, and at $\sigma_a=1\ \mathrm{g\,m^{-2}}$ a mass of $9.4\times10^{17}$ kg. To collect $10^{15}$ W at $\eta=0.3$ and flux $1361\ \mathrm{W\,m^{-2}}$ needs $2.4\times10^{12}\ \mathrm{m^2}$ and $2.4\times10^{9}$ kg. The binding energy of a uniform sphere is $\tfrac35GM^2/R$; for Mercury ($M=3.301\times10^{23}$ kg, $R=2.440\times10^6$ m) it is $1.79\times10^{30}$ J, equal to $4.67\times10^3$ s of solar luminosity; a disassembly over $10^2$ years at 100 per cent efficiency requires $5.7\times10^{20}$ W, a fraction $1.5\times10^{-6}$ of the output.

For the slingshot, let a payload pass the Sun at perihelion radius $r_p$ with speed $v_p$ in the heliocentric frame and make an impulsive burn $\Delta v$ in the direction of motion.

\begin{proposition}
The hyperbolic excess speed after the burn satisfies $v_\infty^2=(v_p+\Delta v)^2-v_{\mathrm{esc}}^2(r_p)=2v_p\Delta v+\Delta v^2-(v_{\mathrm{esc}}^2-v_p^2)$, where $v_{\mathrm{esc}}^2=2GM/r_p$. The specific energy gained is $v_p\Delta v+\tfrac12\Delta v^2$ and is proportional to $v_p$, so for fixed $\Delta v$ the gain is greater at small $r_p$. The velocity increment to reach a perihelion $r_p$ from a circular orbit of radius $r_1$ is $\Delta v_{\mathrm{dive}}=\sqrt{GM/r_1}-\sqrt{GM(2/r_1-1/a)}$ with $a=(r_p+r_1)/2$.
\end{proposition}

\begin{proof}
Energy conservation gives $\tfrac12v_\infty^2=\tfrac12(v_p+\Delta v)^2-GM/r_p$. The dive formula applies the vis-viva equation at the aphelion of the transfer ellipse.
\end{proof}

With $GM_\odot=1.327\times10^{20}\ \mathrm{m^3\,s^{-2}}$, $r_p=10R_\odot=6.957\times10^9$ m, $r_1=1$ AU: $v_{\mathrm{esc}}(r_p)=195.3$ km/s, perihelion speed on the dive ellipse $v_p=190.9$ km/s, $\Delta v_{\mathrm{dive}}=20.9$ km/s. A perihelion burn of $5$, $10$, $15$, $20$ km/s gives $v_\infty=15.5$, $47.2$, $65.3$, $79.6$ km/s at a total velocity increment of $25.9$, $30.9$, $35.9$, $40.9$ km/s. A direct departure from Earth orbit to the same $v_\infty$ requires heliocentric speed $\sqrt{v_\infty^2+2GM/r_1}$ at $1$ AU and thus increments of $15.1$, $33.4$, $47.9$, $60.3$ km/s; the ratio of direct to slingshot cost is $0.58$, $1.08$, $1.33$, $1.47$, with break-even near $v_\infty=40$ km/s. The thermal load at $r_p=10R_\odot$ is $S(1\ \mathrm{AU})(r_1/r_p)^2=6.29\times10^5\ \mathrm{W\,m^{-2}}$ and the blackbody equilibrium temperature of a surface radiating from one side is $(S/\sigma_{SB})^{1/4}=1825$ K.

For the Mercury uses, let $r_m=5.79\times10^{10}$ m be Mercury's orbital radius and $v_m=\sqrt{GM_\odot/r_m}=47.87$ km/s its circular speed.

\begin{proposition}
Let $E_a$ be the embodied energy per unit area of a collector including its machinery, $\eta\Phi$ the useful power per unit area at flux $\Phi$, and $t_p=E_a/(\eta\Phi)$ the energy payback time. If the installed area is reinvested in its own replication, it grows as $A(t)=A_0e^{t/t_p}$, with doubling time $t_d=t_p\ln2$; the number of doublings to reach $A_f$ is $\log_2(A_f/A_0)$. A transport of stored energy $e_b$ per unit mass over a velocity increment $\Delta v$ is energetically self-supporting by propulsion only if $e_b>\tfrac12\Delta v^2$.
\end{proposition}

\begin{proof}
With reinvestment of the whole output, $dA/dt=A\,\eta\Phi/E_a=A/t_p$, which gives the exponential and, setting $A=2A_0$, the doubling time $t_p\ln2$. The condition on transport compares the stored energy with the minimum kinetic energy of the increment.
\end{proof}

Numerically, with $\eta=0.3$, the useful power is $0.408\ \mathrm{kW\,m^{-2}}$ at $1$ AU and $2.72\ \mathrm{kW\,m^{-2}}$ at Mercury. A film of $1\ \mathrm{g\,m^{-2}}$ at $200\ \mathrm{MJ\,kg^{-1}}$ has $E_a=0.2\ \mathrm{MJ\,m^{-2}}$ and $t_p=490$ s at $1$ AU and $73$ s at Mercury. With $100\ \mathrm{kg\,m^{-2}}$ of machinery at the same specific embodied energy, $E_a=2.0\times10^{10}\ \mathrm{J\,m^{-2}}$, $t_p=1.55$ yr and $0.23$ yr, and $t_d=1.08$ yr and $0.16$ yr, respectively; and $\log_2(2.4\times10^{12}/10^{6})=21.2$ doublings. The time to a petawatt swarm is dominated by the machinery term, and $t_d\ge t_p\ln2$ bounds any scheme whose output is the sole source of its own growth. The embodied-energy figures are stipulated inputs of the book and not measurements.

For the dive from Mercury's orbit to $r_p=10R_\odot$, $\Delta v_{\mathrm{dive}}=25.70$ km/s and $v_p=184.6$ km/s by the formula of the preceding proposition with $r_1=r_m$; for $\Delta v=15,20,25$ km/s the excess speed is $40.9,60.8,75.9$ km/s at totals $40.7,45.7,50.7$ km/s, and a direct departure from $r_m$ to the same $v_\infty$ requires $\sqrt{v_\infty^2+2v_m^2}-v_m$, namely $31.2,43.1,53.8$ km/s, so the slingshot breaks even near $v_\infty=60$ km/s. The Hohmann transfer from Mercury's orbit to the Earth's has $\Delta v=9.61+7.53=17.14$ km/s in $105$ days, with $\tfrac12\Delta v^2=147\ \mathrm{MJ\,kg^{-1}}$ against $e_b=0.9\ \mathrm{MJ\,kg^{-1}}$ for a cell of $250\ \mathrm{Wh\,kg^{-1}}$, a ratio of $163$; the escape energy $\tfrac12v_{\mathrm{esc}}^2=9.0\ \mathrm{MJ\,kg^{-1}}$ for $v_{\mathrm{esc}}=4.25$ km/s. The fraction of Mercury used by a one-petawatt swarm is $2.4\times10^{9}/3.30\times10^{23}=7\times10^{-15}$.
